\documentclass[12pt,a4paper]{article}

\usepackage{iftex}
\ifTUTeX % tectonic / xelatex / lualatex
  \usepackage{fontspec}
\else % pdflatex
  \usepackage[utf8]{inputenc}
  \usepackage[T1]{fontenc}
  \usepackage{lmodern}
\fi
\usepackage[polish]{babel}
\usepackage[a4paper,margin=2.5cm]{geometry}
\usepackage{graphicx}
\usepackage{booktabs}
\usepackage{amsmath}
\setlength{\emergencystretch}{2em}
\usepackage{enumitem}
\usepackage{indentfirst}
\usepackage{tikz}
\usetikzlibrary{positioning,fit,backgrounds}
\tikzset{
  wifiicon/.pic={\fill (0,0) circle (0.05);
    \foreach \r in {0.14,0.25,0.36} \draw[line width=1.1pt] (50:\r) arc (50:130:\r);}}
\colorlet{wifibg}{cyan!25}
\usepackage[hidelinks]{hyperref}


% Dane grupy - uzupełnić
\newcommand{\autorA}{Adam Korwin 284455}
\newcommand{\autorB}{Michał Hałas 277462}

\begin{document}

% ---------------- Strona tytułowa ----------------
\pagenumbering{Alph}
\begin{titlepage}
  \centering
  \includegraphics[height=3.5cm]{pwr_logo.png}\par\vspace{0.6cm}
  {\Large\textsc{Politechnika Wrocławska}\par}
  \vspace{0.2cm}
  {\large Wydział Informatyki i Telekomunikacji\par}
  \vspace{2.5cm}
  {\large Projekt z przedmiotu\par}
  \vspace{0.3cm}
  {\LARGE\bfseries Technologie Sieciowe\par}
  \vspace{1.2cm}
  {\Large Projekt lokalnej sieci komputerowej\\ dla wydawnictwa\par}
  \vspace{0.6cm}
  \vspace{2.5cm}
  \begin{flushleft}
    \textbf{Skład grupy:}
    \begin{enumerate}[leftmargin=2em]
      \item \autorA
      \item \autorB
    \end{enumerate}
  \end{flushleft}
  \vfill
  {\large Wrocław, \today\par}
\end{titlepage}
\pagenumbering{arabic}

\tableofcontents
\newpage

% ---------------- 1. Wstęp ----------------
\section{Wstęp}

\subsection{Charakterystyka przedsiębiorstwa}

Przedmiotem projektu jest sieć komputerowa wydawnictwa z siedzibą we Wrocławiu.
Wydawnictwo prowadzi pełny cykl wydawniczy: od przyjęcia materiałów od autorów
i klientów, przez opracowanie redakcyjne i korektę tekstów, aż po przygotowanie
do druku i druk. Firma zatrudnia 310 pracowników pracujących przy stanowiskach
komputerowych, zorganizowanych w trzy grupy robocze (działy):

\begin{itemize}
  \item \textbf{Drukarnia} (103 stanowiska) -- dział odpowiedzialny za
  przygotowanie publikacji do druku (skład, przygotowanie plików produkcyjnych)
  oraz realizację zleceń drukarskich;
  \item \textbf{Korektorzy} (85 stanowisk) -- dział zajmujący się redakcją
  i korektą tekstów przed ich skierowaniem do druku;
  \item \textbf{Konsultanci} (122 stanowiska) -- dział obsługi autorów
  i klientów: przyjmowanie zamówień, konsultacje oraz kontakt z kontrahentami.
\end{itemize}

Przedsiębiorstwo zajmuje dwa budynki: w Budynku~1 piętra 1--2, a w
Budynku~2 piętra 1--3. Pracownicy wszystkich działów są rozmieszczeni
w obu budynkach. Budynki są oddalone od siebie o 204~m i połączone
światłowodem wielomodowym.
W budynkach zainstalowano okablowanie strukturalne kategorii~6 oraz szafy
teleinformatyczne: główny punkt dystrybucyjny (MDF) na 3.~piętrze Budynku~2
oraz dwa pośrednie punkty dystrybucyjne (IDF1 na 2.~piętrze Budynku~2,
IDF2 na 2.~piętrze Budynku~1). Na 1.~piętrze każdego z budynków pracownicy
korzystają również z urządzeń bezprzewodowych. Rozmieszczenie stanowisk,
urządzeń i punktów dystrybucyjnych przedstawiono na rysunku~\ref{fig:rozmieszczenie}.

\begin{figure}[htbp]
\centering
\begin{tikzpicture}[font=\scriptsize,
  floor/.style={draw, anchor=south west, minimum width=6.9cm, minimum height=2cm,
    text width=6.6cm, align=left, inner sep=4pt},
  wifi/.style={fill=wifibg},
  rack/.style={draw, fill=black!75, text=white, font=\scriptsize\bfseries,
    minimum width=1.1cm, minimum height=0.5cm, inner sep=1pt},
  fiber/.style={very thick, orange!80!black},
  copper/.style={very thick, green!50!black},
  serves/.style={densely dotted, thick, ->, >=stealth}]

  % Budynek 1 (2 piętra), punkt dystrybucyjny w prawym górnym rogu
  \node[floor, wifi] (b1p1) at (0,0)
    {\textbf{Piętro 1} -- 99 stanowisk\\
     Drukarnia: 17, Korektorzy: 43, Konsultanci: 39\\
     Drukarki: 1, AP Wi-Fi: 2, urz. bezprz.: 15};
  \node[floor] (b1p2) at (0,2)
    {\textbf{Piętro 2} -- 64 stanowiska\\
     Drukarnia: 22, Korektorzy: 17, Konsultanci: 25\\
     Drukarki: 3};
  \begin{scope}[on background layer]
    \node[draw, very thick, rounded corners, fill=black!4, inner sep=5pt, fit=(b1p1)(b1p2)] (bud1) {};
  \end{scope}
  \node[font=\small\bfseries, below=2pt of bud1] {Budynek 1};

  % Budynek 2 (3 piętra), punkty dystrybucyjne w prawym górnym rogu
  \node[floor, wifi] (b2p1) at (8.6,0)
    {\textbf{Piętro 1} -- 8 stanowisk\\
     Drukarnia: 5, Korektorzy: 2, Konsultanci: 1\\
     Drukarki: 2, AP Wi-Fi: 1, urz. bezprz.: 12};
  \node[floor] (b2p2) at (8.6,2)
    {\textbf{Piętro 2} -- 77 stanowisk\\
     Drukarnia: 24, Korektorzy: 18, Konsultanci: 35\\
     Drukarki: 2};
  \node[floor] (b2p3) at (8.6,4)
    {\textbf{Piętro 3} -- 62 stanowiska\\
     Drukarnia: 35, Korektorzy: 5, Konsultanci: 22\\
     Drukarki: 2};
  \begin{scope}[on background layer]
    \node[draw, very thick, rounded corners, fill=black!4, inner sep=5pt, fit=(b2p1)(b2p3)] (bud2) {};
  \end{scope}
  \node[font=\small\bfseries, below=2pt of bud2] {Budynek 2};

  % Odległość między budynkami
  \draw[<->, >=stealth, thick] (bud1.east |- 0,1) -- (bud2.west |- 0,1)
    node[midway, above, font=\scriptsize\bfseries] {204 m};

  % Wi-Fi
  \pic[cyan!40!black] at ([shift={(-0.45,0.3)}]b1p1.south east) {wifiicon};
  \pic[cyan!40!black] at ([shift={(-0.45,0.3)}]b2p1.south east) {wifiicon};

  \node[rack, anchor=north east] (idf2) at ([shift={(-0.12,-0.12)}]b1p2.north east) {IDF2};
  \node[rack, anchor=north east] (idf1) at ([shift={(-0.12,-0.12)}]b2p2.north east) {IDF1};
  \node[rack, anchor=north east] (mdf)  at ([shift={(-0.12,-0.12)}]b2p3.north east) {MDF};

  \draw[fiber] (idf2.east) -- ++(0.9,0) |- ([yshift=0.35cm]b2p3.north -| mdf) -- (mdf.north);
  \draw[copper] (idf1.north) -- (mdf.south);
  \draw[serves] ([xshift=0.35cm]idf2.south) -- ([xshift=-0.32cm]b1p1.north east);
  \draw[serves] ([xshift=0.35cm]idf1.south) -- ([xshift=-0.32cm]b2p1.north east);
\end{tikzpicture}

\vspace{6pt}
{\scriptsize
\tikz\draw[very thick, orange!80!black] (0,0) -- (0.6,0); światłowód wielomodowy\quad
\tikz\draw[very thick, green!50!black] (0,0) -- (0.6,0); skrętka kat.~6\quad
\tikz\draw[densely dotted, thick, ->, >=stealth] (0,0) -- (0.6,0); piętro obsługiwane przez IDF\quad
\tikz\draw[fill=wifibg] (0,0) rectangle (0.4,0.22); \tikz[baseline=-1pt]\pic[cyan!40!black,scale=0.8]{wifiicon}; piętro z siecią Wi-Fi\par}
\caption{Rozmieszczenie stanowisk, urządzeń i punktów dystrybucyjnych w budynkach wydawnictwa}
\label{fig:rozmieszczenie}
\end{figure}

Profil działalności wydawnictwa przekłada się bezpośrednio na charakter ruchu
w sieci. Praca na tekstach i plikach przygotowanych do druku wymaga intensywnej
wymiany danych z serwerami lokalnymi oraz z drukarkami sieciowymi, w których
dominuje ruch wysyłany przez użytkowników (zapis dokumentów, zlecenia wydruku).
Firma utrzymuje także własne serwery dostępne z Internetu: serwer WWW
(serwis internetowy wydawnictwa) oraz serwer FTP, służący do wymiany plików
z autorami i klientami. Pracownicy korzystają ponadto z aplikacji
internetowych: przeglądarki, wideokonferencji, telefonii VoIP, klienta FTP,
komunikatora, usług pracy w chmurze oraz poczty elektronicznej. Ciągłość pracy
sieci jest dla firmy istotna, ponieważ jej przerwa zatrzymuje zarówno
produkcję wydawniczą, jak i obsługę klientów.

\subsection{Cel projektu}

Celem projektu jest zaprojektowanie lokalnej sieci komputerowej wydawnictwa,
która zapewni prawidłową i niezawodną obsługę wszystkich pracowników, serwerów,
drukarek oraz urządzeń bezprzewodowych, z uwzględnieniem rzeczywistych
wymagań na przepustowość zgłoszonych przez zamawiającego.

Projektowana sieć ma spełniać następujące wymagania zamawiającego:
\begin{itemize}
  \item wykorzystanie technologii z rodziny Ethernet;
  \item separacja ruchu poszczególnych grup roboczych za pomocą sieci VLAN;
  \item dostęp do sieci bezprzewodowej na wskazanym piętrze każdego budynku;
  \item rezerwa portów na przełącznikach w wysokości 20\,\% portów zajętych,
  wynikająca z przewidywanego wzrostu zatrudnienia;
  \item dwa łącza do Internetu (podstawowe i zapasowe) o przepustowości
  adekwatnej do potrzeb firmy, przy czym łącze podstawowe ma gwarantować
  co najmniej 40\,\% średniego przewidywanego przepływu;
  \item zabezpieczenie sieci przed typowymi zagrożeniami, takimi jak awaria
  zasilania, złośliwe oprogramowanie czy próby włamania z Internetu.
\end{itemize}

\subsection{Zakres projektu}

Okablowanie strukturalne, szafy teleinformatyczne oraz urządzenia końcowe
(komputery, serwery, drukarki) są już zainstalowane i nie wchodzą w zakres
prac. Zakres projektu obejmuje:
\begin{enumerate}
  \item inwentaryzację zasobów sprzętowych, aplikacji i zasobów ludzkich,
  \item analizę potrzeb użytkowników, w tym oszacowanie przepływów danych
  w kluczowych punktach sieci,
  \item sformułowanie założeń projektowych,
  \item projekt logiczny sieci wraz z podziałem na sieci VLAN,
  \item dobór urządzeń sieciowych,
  \item projekt adresacji IP oraz konfiguracji urządzeń,
  \item projekt podłączenia do Internetu wraz z wyborem dostawców,
  \item analizę bezpieczeństwa i niezawodności sieci,
  \item kosztorys obejmujący urządzenia oraz koszt łączy internetowych
  w okresie 2~lat.
\end{enumerate}

% ================= 2. Inwentaryzacja =================
\section{Inwentaryzacja zasobów: sprzętu, aplikacji, zasobów ludzkich}

Rozdział przedstawia dane otrzymane od Prowadzącego (zestaw nr~2). Wszystkie
przepływności podano w kb/s; w dalszych obliczeniach przyjęto
$1~\text{Mb/s} = 1000~\text{kb/s}$.

\subsection{Zasoby ludzkie}

W firmie pracuje 310 osób przy stanowiskach komputerowych (jeden pracownik
odpowiada jednemu komputerowi). Pracownicy należą do trzech grup roboczych
i są rozmieszczeni na pięciu piętrach w dwóch budynkach
(tabela~\ref{tab:uzytkownicy}).

\begin{table}[htbp]
\centering
\caption{Liczba użytkowników (komputerów) w grupach roboczych na piętrach}
\label{tab:uzytkownicy}
\begin{tabular}{lcccccr}
\toprule
 & \multicolumn{2}{c}{Budynek 1} & \multicolumn{3}{c}{Budynek 2} & \\
\cmidrule(lr){2-3}\cmidrule(lr){4-6}
Grupa robocza & P.~1 & P.~2 & P.~1 & P.~2 & P.~3 & Razem \\
\midrule
Drukarnia   & 17 & 22 & 5 & 24 & 35 & 103 \\
Korektorzy  & 43 & 17 & 2 & 18 &  5 &  85 \\
Konsultanci & 39 & 25 & 1 & 35 & 22 & 122 \\
\midrule
Razem       & 99 & 64 & 8 & 77 & 62 & \textbf{310} \\
\bottomrule
\end{tabular}
\end{table}

Najliczniejszą grupą są Konsultanci (39\,\% stanowisk), a najwięcej
pracowników jest na 1.~piętrze Budynku~1 (99 stanowisk). Każda grupa jest
rozproszona po wszystkich pięciu piętrach. Z tego wynika, że separacji ruchu
grup roboczych nie da się zrealizować podziałem fizycznym (osobne
przełączniki dla działów), tylko logicznym, za pomocą sieci VLAN obejmujących
oba budynki.

\subsection{Sprzęt}

\paragraph{Urządzenia końcowe na piętrach.} Oprócz komputerów na piętrach
znajdują się drukarki sieciowe i punkty dostępowe Wi-Fi, z których korzystają
urządzenia bezprzewodowe (tabela~\ref{tab:sprzet}). Sieć bezprzewodowa ma
działać na 1.~piętrze każdego budynku.

\begin{table}[htbp]
\centering
\caption{Drukarki, punkty dostępowe i urządzenia bezprzewodowe na piętrach}
\label{tab:sprzet}
\begin{tabular}{lcccccr}
\toprule
 & \multicolumn{2}{c}{Budynek 1} & \multicolumn{3}{c}{Budynek 2} & \\
\cmidrule(lr){2-3}\cmidrule(lr){4-6}
Urządzenie & P.~1 & P.~2 & P.~1 & P.~2 & P.~3 & Razem \\
\midrule
Drukarki                    &  1 & 3 &  2 & 2 & 2 & 10 \\
Punkty dostępowe Wi-Fi      &  2 & 0 &  1 & 0 & 0 &  3 \\
Urządzenia bezprzewodowe    & 15 & 0 & 12 & 0 & 0 & 27 \\
\bottomrule
\end{tabular}
\end{table}

\paragraph{Serwery.} Firma posiada dwa serwery lokalne (Serwer1, Serwer2)
oraz dwa serwery udostępniane w Internecie (serwer WWW i serwer FTP). Z~danych
wynika, że z Serwera1 korzystają wszyscy użytkownicy, łącznie z bezprzewodowymi,
a z Serwera2 wyłącznie Korektorzy. Lokalizacja serwerów nie została podana.

\paragraph{Punkty dystrybucyjne i okablowanie.} W budynkach zainstalowano
okablowanie strukturalne kat.~6 oraz trzy szafy teleinformatyczne
(tabela~\ref{tab:pd}). Budynki są oddalone o 204~m i połączone światłowodem
wielomodowym (kategoria włókna, np.\ OM2/OM3, nie została podana). Wewnątrz
budynku między MDF a IDF ułożono skrętkę kat.~6, a odległość między
punktami dystrybucyjnymi nie przekracza 90~m.

\begin{table}[htbp]
\centering
\caption{Punkty dystrybucyjne}
\label{tab:pd}
\begin{tabular}{lll}
\toprule
Oznaczenie & Lokalizacja & Obsługiwane piętra \\
\midrule
MDF  & Budynek 2, piętro 3 & Budynek 2: piętro 3 \\
IDF1 & Budynek 2, piętro 2 & Budynek 2: piętra 1 i 2 \\
IDF2 & Budynek 1, piętro 2 & Budynek 1: piętra 1 i 2 \\
\bottomrule
\end{tabular}
\end{table}

\subsection{Aplikacje i usługi}

\paragraph{Ruch do serwerów lokalnych i drukarek.} Tabela~\ref{tab:lokalne}
podaje przepływności przypadające na jednego użytkownika danej grupy.
Kierunek \emph{down} oznacza ruch do użytkownika, a \emph{up} ruch od
użytkownika do serwera lub drukarki. Wiersz WiFi dotyczy urządzeń
bezprzewodowych.

\begin{table}[htbp]
\centering
\caption{Przepływności do serwerów lokalnych i drukarek na użytkownika (down\,/\,up) [kb/s]}
\label{tab:lokalne}
\begin{tabular}{lccc}
\toprule
Grupa & Serwer1 & Serwer2 & Drukarka \\
\midrule
Drukarnia   & 100 / 700 &   0 / 0   & 10 / 190 \\
Korektorzy  & 200 / 250 & 250 / 750 & 10 / 200 \\
Konsultanci & 350 / 850 &   0 / 0   & 10 / 120 \\
WiFi        &  50 / 100 &   0 / 0   & 10 / 120 \\
\bottomrule
\end{tabular}
\end{table}

W ruchu do serwerów i drukarek przeważa kierunek \emph{up}, czyli zapis plików
i wysyłanie zleceń wydruku. Jest to zgodne z profilem wydawnictwa: pracownicy
zapisują na serwerach przygotowane i poprawione dokumenty.

\paragraph{Aplikacje internetowe.} Tabela~\ref{tab:aplikacje} podaje
przepływności generowane przez aplikacje użytkowników w kierunku Internetu.

\begin{table}[htbp]
\centering
\caption{Przepływności aplikacji internetowych na użytkownika (down\,/\,up) [kb/s]}
\label{tab:aplikacje}
\small
\setlength{\tabcolsep}{4pt}
\begin{tabular}{lccccccc}
\toprule
Grupa & Przegl. & Wideokonf. & VoIP & Klient FTP & Komunik. & Chmura & Poczta \\
\midrule
Drukarnia   & 74 / 10 & 40 / 40 & 20 / 20 & 78 / 11 & 15 / 15 & 54 / 28 & 18 / 26 \\
Korektorzy  & 42 / 10 &  0 / 0  & 20 / 20 & 69 / 20 & 15 / 15 & 47 / 49 &  0 / 0  \\
Konsultanci & 61 / 10 & 40 / 40 &  0 / 0  &  0 / 0  & 15 / 15 &  0 / 0  & 10 / 29 \\
WiFi        &  0 / 0  &  0 / 0  &  0 / 0  & 69 / 17 & 15 / 15 & 44 / 46 &  6 / 17 \\
\bottomrule
\end{tabular}
\end{table}

Wśród aplikacji są usługi czasu rzeczywistego (VoIP, wideokonferencje,
komunikator), które są wrażliwe na opóźnienia i zmienność opóźnień.

\paragraph{Serwery internetowe.} Tabela~\ref{tab:serwery} podaje prognozowany
ruch generowany przez serwery firmy obsługujące użytkowników z Internetu.

\begin{table}[htbp]
\centering
\caption{Prognozowany ruch serwerów internetowych}
\label{tab:serwery}
\begin{tabular}{lccc}
\toprule
Serwer & Do Internetu [kb/s/sesję] & Z Internetu [kb/s/sesję] & Jednoczesne sesje \\
\midrule
WWW & 130 & 25 & 33 \\
FTP & 130 & 50 &  3 \\
\bottomrule
\end{tabular}
\end{table}

% ================= 3. Analiza potrzeb =================
\section{Analiza potrzeb użytkowników -- wymagania zamawiającego}

\subsection{Założenia do obliczeń}

\begin{enumerate}
  \item Podane przepływności uwzględniają przeciążenia i chwilowe wzrosty
  ruchu (zgodnie z wytycznymi), dlatego nie stosuje się dodatkowych
  współczynników. Przyjęto, że wszyscy użytkownicy pracują jednocześnie.
  Suma przepływów jest więc średnim przewidywanym przepływem w danym punkcie
  sieci. Ponieważ dane od zamawiającego zawierają już zapas na chwilowe wzrosty
  ruchu, ta sama wartość jest podstawą doboru przepływności łączy.
  \item Urządzenia bezprzewodowe (27) są użytkownikami dodatkowymi, poza
  310 komputerami, o profilu ruchu z wiersza WiFi. Łącznie sieć obsługuje
  $103 + 85 + 122 + 27 = 337$ użytkowników.
  \item Serwery lokalne (Serwer1, Serwer2) oraz serwery internetowe (WWW, FTP)
  znajdują się w MDF (Budynek~2, piętro~3), w którym zbiega się okablowanie
  obu budynków. Serwery lokalne są podłączone do przełącznika rdzeniowego,
  a serwery WWW i FTP do strefy DMZ zapory sieciowej.
  \item Użytkownik drukuje na drukarkach swojego piętra, a ruch rozkłada się
  równo na drukarki piętra.
  \item Drukarki są wspólne dla wszystkich grup, więc będą w osobnej sieci
  VLAN, a routing między VLAN-ami odbywa się w MDF. Ruch od użytkownika do
  drukarki w budynku obsługiwanym przez IDF przechodzi więc łączem
  szkieletowym dwukrotnie: do MDF i z powrotem. Dlatego w obliczeniach łączy
  szkieletowych ruch drukarek (down + up) występuje w obu kierunkach. Wynika to
  bezpośrednio z umieszczenia routingu w MDF. Gdyby drukarki znalazły się
  w VLAN-ach działów albo routing odbywał się w IDF, ruch do drukarek
  pozostałby w obrębie budynku.
\end{enumerate}

Przepływ w danym punkcie sieci obliczono jako
\begin{equation}
  T = \sum_{g} n_g \cdot t_g ,
  \label{eq:suma}
\end{equation}
gdzie $n_g$ to liczba użytkowników grupy $g$ korzystających z danego punktu
sieci, a $t_g$ to przepływność na jednego użytkownika tej grupy.

\subsection{Ruch do serwerów lokalnych}

Łącznie z Serwera1 korzysta 337 użytkowników (103~D, 85~K, 122~S, 27~WiFi).
\begin{align*}
T_{\text{S1,up}}   &= 103\cdot700 + 85\cdot250 + 122\cdot850 + 27\cdot100 \\
                    &= 72\,100 + 21\,250 + 103\,700 + 2\,700 = 199\,750~\text{kb/s},\\
T_{\text{S1,down}} &= 103\cdot100 + 85\cdot200 + 122\cdot350 + 27\cdot50 \\
                    &= 10\,300 + 17\,000 + 42\,700 + 1\,350 = 71\,350~\text{kb/s}.
\end{align*}
Z Serwera2 korzystają tylko Korektorzy (85 osób):
\begin{align*}
T_{\text{S2,up}}   &= 85\cdot750 = 63\,750~\text{kb/s}, &
T_{\text{S2,down}} &= 85\cdot250 = 21\,250~\text{kb/s}.
\end{align*}
Łącze Serwera1 musi więc przyjmować ok.\ 199{,}8~Mb/s i wysyłać 71{,}4~Mb/s,
a łącze Serwera2 przyjmować 63{,}8~Mb/s i wysyłać 21{,}3~Mb/s.

\subsection{Ruch do drukarek}

Łączny ruch do wszystkich drukarek wynosi
\begin{align*}
T_{\text{Dr,up}}   &= 103\cdot190 + 85\cdot200 + 122\cdot120 + 27\cdot120 \\
                    &= 19\,570 + 17\,000 + 14\,640 + 3\,240 = 54\,450~\text{kb/s},\\
T_{\text{Dr,down}} &= 337\cdot10 = 3\,370~\text{kb/s}.
\end{align*}
Tabela~\ref{tab:drukarki} pokazuje obciążenie łącza pojedynczej drukarki na
każdym piętrze. Ruch piętra dzieli się przez liczbę drukarek na piętrze
($T_\text{down}$ to 10~kb/s na każdego użytkownika piętra):
\begin{align*}
\text{B1/P1:}~T_\text{up} &= 17\cdot190 + 43\cdot200 + 39\cdot120 + 15\cdot120 = 18\,310,
  & T_\text{down} &= 114\cdot10 = 1\,140;\\
\text{B1/P2:}~T_\text{up} &= 22\cdot190 + 17\cdot200 + 25\cdot120 = 10\,580,
  & T_\text{down} &= 64\cdot10 = 640;\\
\text{B2/P1:}~T_\text{up} &= 5\cdot190 + 2\cdot200 + 1\cdot120 + 12\cdot120 = 2\,910,
  & T_\text{down} &= 20\cdot10 = 200;\\
\text{B2/P2:}~T_\text{up} &= 24\cdot190 + 18\cdot200 + 35\cdot120 = 12\,360,
  & T_\text{down} &= 77\cdot10 = 770;\\
\text{B2/P3:}~T_\text{up} &= 35\cdot190 + 5\cdot200 + 22\cdot120 = 10\,290,
  & T_\text{down} &= 62\cdot10 = 620
\end{align*}
(wartości w kb/s). Na przykład na 2.~piętrze Budynku~1 są 3 drukarki, więc na
jedną przypada $10\,580/3 = 3\,526{,}7$~kb/s up i $640/3 = 213{,}3$~kb/s down.

\begin{table}[htbp]
\centering
\caption{Obciążenie łącza jednej drukarki na piętrach}
\label{tab:drukarki}
\small
\begin{tabular}{lccrrr}
\toprule
Piętro & Użytk. & Liczba & Up piętra & Up na & Down na \\
 & D/K/S/WiFi & drukarek & [kb/s] & drukarkę [kb/s] & drukarkę [kb/s] \\
\midrule
B1, P.~1 & 17/43/39/15 & 1 & 18\,310 & 18\,310{,}0 & 1\,140{,}0 \\
B1, P.~2 & 22/17/25/0  & 3 & 10\,580 &  3\,526{,}7 &   213{,}3 \\
B2, P.~1 & 5/2/1/12    & 2 &  2\,910 &  1\,455{,}0 &   100{,}0 \\
B2, P.~2 & 24/18/35/0  & 2 & 12\,360 &  6\,180{,}0 &   385{,}0 \\
B2, P.~3 & 35/5/22/0   & 2 & 10\,290 &  5\,145{,}0 &   310{,}0 \\
\midrule
Razem    &             & 10 & 54\,450 & & \\
\bottomrule
\end{tabular}
\end{table}

Najbardziej obciążona jest jedyna drukarka na 1.~piętrze Budynku~1
(ok.\ 18{,}3~Mb/s). Łącze 100~Mb/s wystarcza dla każdej drukarki.

Pojedyncze stanowisko generuje najwyżej $t_\text{do MDF} = 1\,324$~kb/s
(Korektor, tabela~\ref{tab:naosobe}), więc port 1~Gb/s jest obciążony
w znikomym stopniu.

\subsection{Ruch do Internetu}

\paragraph{Użytkownicy.} Przepływność na jednego użytkownika to suma
przepływności wszystkich aplikacji z tabeli~\ref{tab:aplikacje}:
\begin{align*}
\text{Drukarnia:}   \quad & t_\text{down} = 74+40+20+78+15+54+18 = 299, \\
                          & t_\text{up} = 10+40+20+11+15+28+26 = 150~\text{kb/s},\\
\text{Korektorzy:}  \quad & t_\text{down} = 42+0+20+69+15+47+0 = 193, \\
                          & t_\text{up} = 10+0+20+20+15+49+0 = 114~\text{kb/s},\\
\text{Konsultanci:} \quad & t_\text{down} = 61+40+0+0+15+0+10 = 126, \\
                          & t_\text{up} = 10+40+0+0+15+0+29 = 94~\text{kb/s},\\
\text{WiFi:}        \quad & t_\text{down} = 0+0+0+69+15+44+6 = 134, \\
                          & t_\text{up} = 0+0+0+17+15+46+17 = 95~\text{kb/s}.
\end{align*}
Ruch wszystkich użytkowników:
\begin{align*}
T_\text{U,down} &= 103\cdot299 + 85\cdot193 + 122\cdot126 + 27\cdot134 \\
                 &= 30\,797 + 16\,405 + 15\,372 + 3\,618 = 66\,192~\text{kb/s},\\
T_\text{U,up}   &= 103\cdot150 + 85\cdot114 + 122\cdot94 + 27\cdot95 \\
                 &= 15\,450 + 9\,690 + 11\,468 + 2\,565 = 39\,173~\text{kb/s}.
\end{align*}

\paragraph{Serwery internetowe.} Ruch serwerów to iloczyn przepływności
na sesję i liczby jednoczesnych sesji:
\begin{align*}
T_\text{Srv,do Int.} &= 33\cdot130 + 3\cdot130 = 4\,290 + 390 = 4\,680~\text{kb/s},\\
T_\text{Srv,z Int.}  &= 33\cdot25 + 3\cdot50 = 825 + 150 = 975~\text{kb/s}.
\end{align*}
Łącze serwera WWW obsługuje 4{,}29~Mb/s wysyłanych i 0{,}83~Mb/s
odbieranych, a łącze serwera FTP 0{,}39~Mb/s i 0{,}15~Mb/s.

\paragraph{Łącze internetowe.} Kierunek do firmy (\emph{download}) tworzą
dane pobierane przez użytkowników i zapytania do serwerów firmy, a kierunek
do Internetu (\emph{upload}) dane wysyłane przez użytkowników i odpowiedzi
serwerów:
\begin{align*}
T_\text{Int,down} &= 66\,192 + 975    = 67\,167~\text{kb/s} \approx 67{,}2~\text{Mb/s},\\
T_\text{Int,up}   &= 39\,173 + 4\,680 = 43\,853~\text{kb/s} \approx 43{,}9~\text{Mb/s}.
\end{align*}
Łącze podstawowe ma gwarantować co najmniej 40\,\% średniego przewidywanego
przepływu:
\begin{align*}
G_\text{down} &= 0{,}4\cdot67\,167 = 26\,866{,}8~\text{kb/s} \approx 26{,}9~\text{Mb/s},\\
G_\text{up}   &= 0{,}4\cdot43\,853 = 17\,541{,}2~\text{kb/s} \approx 17{,}5~\text{Mb/s}.
\end{align*}

\subsection{Łącza szkieletowe}

Przez łącze IDF--MDF przechodzi cały ruch użytkowników danego IDF do serwerów,
do Internetu oraz, zgodnie z założeniem~5, do drukarek w obu kierunkach. Na
jednego użytkownika przypada więc:
\begin{align*}
t_\text{do MDF} &= t_\text{S1,up} + t_\text{S2,up} + t_\text{Int,up} + (t_\text{Dr,down} + t_\text{Dr,up}),\\
t_\text{z MDF}  &= t_\text{S1,down} + t_\text{S2,down} + t_\text{Int,down} + (t_\text{Dr,down} + t_\text{Dr,up}).
\end{align*}

\begin{table}[htbp]
\centering
\caption{Łączny ruch na jednego użytkownika na łączu szkieletowym [kb/s]}
\label{tab:naosobe}
\begin{tabular}{lll}
\toprule
Grupa & Do MDF & Z MDF \\
\midrule
Drukarnia   & $700 + 0 + 150 + 200 = 1050$   & $100 + 0 + 299 + 200 = 599$ \\
Korektorzy  & $250 + 750 + 114 + 210 = 1324$ & $200 + 250 + 193 + 210 = 853$ \\
Konsultanci & $850 + 0 + 94 + 130 = 1074$    & $350 + 0 + 126 + 130 = 606$ \\
WiFi        & $100 + 0 + 95 + 130 = 325$     & $50 + 0 + 134 + 130 = 314$ \\
\bottomrule
\end{tabular}
\end{table}

\paragraph{Łącze IDF2--MDF (między budynkami, światłowód).} IDF2 obsługuje
cały Budynek~1: $17+22 = 39$~D, $43+17 = 60$~K, $39+25 = 64$~S oraz 15~WiFi.
\begin{align*}
T_\text{do MDF} &= 39\cdot1050 + 60\cdot1324 + 64\cdot1074 + 15\cdot325 \\
                 &= 40\,950 + 79\,440 + 68\,736 + 4\,875 = 194\,001~\text{kb/s},\\
T_\text{z MDF}  &= 39\cdot599 + 60\cdot853 + 64\cdot606 + 15\cdot314 \\
                 &= 23\,361 + 51\,180 + 38\,784 + 4\,710 = 118\,035~\text{kb/s}.
\end{align*}

\paragraph{Łącze IDF1--MDF (Budynek 2, skrętka).} IDF1 obsługuje piętra 1 i 2
Budynku~2: $5+24 = 29$~D, $2+18 = 20$~K, $1+35 = 36$~S oraz 12~WiFi.
\begin{align*}
T_\text{do MDF} &= 29\cdot1050 + 20\cdot1324 + 36\cdot1074 + 12\cdot325 \\
                 &= 30\,450 + 26\,480 + 38\,664 + 3\,900 = 99\,494~\text{kb/s},\\
T_\text{z MDF}  &= 29\cdot599 + 20\cdot853 + 36\cdot606 + 12\cdot314 \\
                 &= 17\,371 + 17\,060 + 21\,816 + 3\,768 = 60\,015~\text{kb/s}.
\end{align*}

\paragraph{Przełączniki dostępowe w MDF.} Piętro~3 Budynku~2 (35~D, 5~K,
22~S) łączy się z przełącznikiem rdzeniowym wewnątrz szafy MDF:
\begin{align*}
T_\text{do rdzenia} &= 35\cdot1050 + 5\cdot1324 + 22\cdot1074 \\
                     &= 36\,750 + 6\,620 + 23\,628 = 66\,998~\text{kb/s},\\
T_\text{z rdzenia}  &= 35\cdot599 + 5\cdot853 + 22\cdot606 \\
                     &= 20\,965 + 4\,265 + 13\,332 = 38\,562~\text{kb/s}.
\end{align*}

Kontrola: suma ruchu do rdzenia z trzech punktów dystrybucyjnych,
$194\,001 + 99\,494 + 66\,998 = 360\,493$~kb/s, jest równa sumie przepływów
od wszystkich użytkowników:
\[ 199\,750 + 63\,750 + 39\,173 + (54\,450 + 3\,370) = 360\,493~\text{kb/s}. \]

\subsection{Sieć bezprzewodowa}

Na jedno urządzenie bezprzewodowe przypada
$t_\text{down} = 50 + 0 + 10 + 134 = 194$~kb/s oraz
$t_\text{up} = 100 + 0 + 120 + 95 = 315$~kb/s.
\begin{itemize}
  \item Budynek~1, piętro~1: 15 urządzeń, 2 punkty dostępowe, czyli 7,5
  urządzenia na punkt: $7{,}5\cdot194 = 1\,455$~kb/s down
  i $7{,}5\cdot315 = 2\,362{,}5$~kb/s up.
  \item Budynek~2, piętro~1: 12 urządzeń, 1 punkt dostępowy:
  $12\cdot194 = 2\,328$~kb/s down i $12\cdot315 = 3\,780$~kb/s up.
\end{itemize}
Wymagana przepływność jest niewielka (poniżej 4~Mb/s na punkt dostępowy).
O doborze punktów dostępowych decyduje więc zasięg i liczba jednocześnie
obsługiwanych urządzeń, a nie przepływność.

\subsection{Liczba portów}

Zgodnie z wymaganiem przełączniki mają mieć 20\,\% portów zapasowych
w stosunku do portów zajętych. Liczbę portów dla urządzeń końcowych
w każdym punkcie dystrybucyjnym podano w tabeli~\ref{tab:porty} (wynik
zaokrąglono w górę). Serwery nie są podłączone do przełączników dostępowych,
dlatego ich porty policzono osobno: Serwer1 i Serwer2 po 2~porty (agregacja
LACP) na przełączniku rdzeniowym oraz serwery WWW i FTP po 1~porcie w strefie
DMZ zapory. Daje to $6\cdot1{,}2 = 7{,}2 \to 8$ portów. Porty łączy
szkieletowych oraz porty do routerów i zapory zostaną policzone na etapie
doboru urządzeń.

\begin{table}[htbp]
\centering
\caption{Liczba portów dostępowych w punktach dystrybucyjnych}
\label{tab:porty}
\small
\begin{tabular}{lcccl}
\toprule
Punkt & Komputery & Drukarki & AP & Zajęte $\to$ z rezerwą 20\,\% \\
\midrule
IDF2 (B1, P.~1--2) & 163 & 4 & 2  & $169 \cdot 1{,}2 = 202{,}8 \to 203$ \\
IDF1 (B2, P.~1--2) &  85 & 4 & 1  & $90 \cdot 1{,}2 = 108$ \\
MDF  (B2, P.~3)    &  62 & 2 & -- & $64 \cdot 1{,}2 = 76{,}8 \to 77$ \\
\midrule
Razem              & 310 & 10 & 3 & $323 \to 388$ \\
\bottomrule
\end{tabular}
\end{table}

\subsection{Łącze do zapory sieciowej i strefa DMZ}

Serwery WWW i FTP są podłączone do strefy DMZ zapory, więc łączem przełącznik
rdzeniowy--zapora płynie tylko ruch użytkowników do Internetu:
$T_\text{U,down} = 66{,}2$~Mb/s (do rdzenia) i $T_\text{U,up} = 39{,}2$~Mb/s
(do zapory). Łączem zapora--DMZ płynie ruch serwerów internetowych:
4{,}68~Mb/s z DMZ i 0{,}98~Mb/s do DMZ. Przez zaporę przechodzi łącznie cały
ruch łącza internetowego, czyli 67{,}2/43{,}9~Mb/s.

Przełącznik rdzeniowy obsługuje routing między VLAN-ami. Przechodzi przez
niego cały ruch od użytkowników, równy sumie z kontroli w punkcie~3.5
(360{,}5~Mb/s), oraz ruch w przeciwnym kierunku.

\subsection{Rezerwa na wzrost liczby użytkowników}

Zamawiający przewiduje wzrost liczby pracowników o 20\,\%. Zakładając, że nowi
użytkownicy mają ten sam profil ruchu, przepływy generowane przez użytkowników
rosną 1,2 raza, a ruch serwerów internetowych pozostaje bez zmian.
Na przykład dla łącza internetowego:
\begin{align*}
T_\text{Int,down} &= 1{,}2\cdot66\,192 + 975 = 79\,430{,}4 + 975 = 80\,405{,}4~\text{kb/s} \approx 80{,}4~\text{Mb/s},\\
T_\text{Int,up}   &= 1{,}2\cdot39\,173 + 4\,680 = 47\,007{,}6 + 4\,680 = 51\,687{,}6~\text{kb/s} \approx 51{,}7~\text{Mb/s}.
\end{align*}

\begin{table}[htbp]
\centering
\caption{Przepływy po wzroście liczby użytkowników o 20\,\% [Mb/s]}
\label{tab:wzrost}
\begin{tabular}{lcc}
\toprule
Punkt sieci & Obecnie & Po wzroście o 20\,\% \\
\midrule
Łącze IDF2--MDF (do / z MDF)     & 194{,}0 / 118{,}0 & 232{,}8 / 141{,}6 \\
Łącze IDF1--MDF (do / z MDF)     & 99{,}5 / 60{,}0   & 119{,}4 / 72{,}0 \\
Serwer1 (do / od serwera)        & 199{,}8 / 71{,}4  & 239{,}7 / 85{,}6 \\
Łącze do Internetu (down / up)   & 67{,}2 / 43{,}9   & 80{,}4 / 51{,}7 \\
Gwarancja 40\,\% (down / up)     & 26{,}9 / 17{,}5   & 32{,}2 / 20{,}7 \\
\bottomrule
\end{tabular}
\end{table}

\subsection{Podsumowanie wymagań na przepływność}

\begin{table}[htbp]
\centering
\caption{Wymagana przepływność w kluczowych punktach sieci}
\label{tab:podsumowanie}
\begin{tabular}{lrr}
\toprule
Punkt sieci & Kierunek 1 [Mb/s] & Kierunek 2 [Mb/s] \\
\midrule
Łącze IDF2--MDF (między budynkami) & do MDF: 194{,}0 & z MDF: 118{,}0 \\
Łącze IDF1--MDF                    & do MDF: 99{,}5  & z MDF: 60{,}0 \\
Przełączniki piętra 3 -- rdzeń (MDF) & do rdzenia: 67{,}0 & z rdzenia: 38{,}6 \\
Serwer1                            & do serwera: 199{,}8 & od serwera: 71{,}4 \\
Serwer2                            & do serwera: 63{,}8  & od serwera: 21{,}3 \\
Serwer WWW                         & do Internetu: 4{,}29 & z Internetu: 0{,}83 \\
Serwer FTP                         & do Internetu: 0{,}39 & z Internetu: 0{,}15 \\
Najbardziej obciążona drukarka     & do drukarki: 18{,}3 & od drukarki: 1{,}1 \\
Rdzeń -- zapora                    & do zapory: 39{,}2 & do rdzenia: 66{,}2 \\
Zapora -- DMZ                      & z DMZ: 4{,}68 & do DMZ: 0{,}98 \\
Łącze do Internetu                 & down: 67{,}2 & up: 43{,}9 \\
Gwarancja łącza podstawowego (40\,\%) & down: 26{,}9 & up: 17{,}5 \\
\bottomrule
\end{tabular}
\end{table}

\subsection{Pozostałe wymagania}

\paragraph{Wymagania funkcjonalne.}
\begin{itemize}
  \item Ruch każdej grupy roboczej ma być odseparowany (VLAN), przy czym grupy
  są rozproszone po obu budynkach.
  \item Dostęp do Serwera2 potrzebują tylko Korektorzy, więc dostęp innych
  grup do tego serwera można zablokować.
  \item Sieć Wi-Fi ma działać na 1.~piętrze obu budynków.
  \item Serwery WWW i FTP muszą być dostępne z Internetu.
  \item Usługi czasu rzeczywistego (VoIP, wideokonferencje) wymagają
  małych opóźnień.
\end{itemize}

\paragraph{Wymagania dotyczące bezpieczeństwa.}
\begin{itemize}
  \item Ochrona sieci wewnętrznej przed atakami z Internetu, w szczególności
  oddzielenie publicznie dostępnych serwerów WWW i FTP od sieci wewnętrznej.
  \item Kontrola dostępu między grupami roboczymi, serwerami i drukarkami.
  \item Zabezpieczenie sieci bezprzewodowej przed dostępem osób
  nieuprawnionych.
  \item Ochrona przed złośliwym oprogramowaniem i utratą danych na serwerach
  (pliki wydawnicze są głównym majątkiem firmy).
  \item Ochrona przed podłączeniem nieautoryzowanych urządzeń do gniazd
  sieciowych.
\end{itemize}

\paragraph{Wymagania dotyczące niezawodności.}
\begin{itemize}
  \item Dwa łącza do Internetu (podstawowe i zapasowe) z automatycznym
  przełączaniem.
  \item Odporność na awarię zasilania urządzeń sieciowych i serwerów.
  \item Odporność na awarię pojedynczego łącza szkieletowego, zwłaszcza
  łącza między budynkami, od którego zależy praca 163 stanowisk.
  \item Możliwość rozbudowy sieci o 20\,\% użytkowników bez wymiany
  urządzeń.
  \item Brak pojedynczych punktów awarii w elementach, od których zależy
  praca całej firmy (przełącznik rdzeniowy, zapora, łącze do Internetu),
  albo świadoma akceptacja tego ryzyka przez zamawiającego.
\end{itemize}

% ================= 4. Założenia projektowe =================
\section{Założenia projektowe}

Na podstawie inwentaryzacji i analizy potrzeb przyjęto następujące założenia
(do wstępnej akceptacji przez zamawiającego).

\paragraph{Technologia i topologia LAN.}
\begin{enumerate}
  \item Technologia Ethernet (IEEE~802.3), topologia gwiazdy hierarchicznej:
  przełącznik rdzeniowy warstwy~3 w MDF, przełączniki dostępowe w MDF, IDF1
  i IDF2.
  \item Stanowiska, drukarki i punkty dostępowe: 1000BASE-T (1~Gb/s) na
  istniejącym okablowaniu kat.~6. Porty punktów dostępowych są zasilane
  przez PoE, więc przełączniki obsługujące piętra 1 (IDF1, IDF2) mają porty
  PoE.
  \item Przełączniki dostępowe w każdym punkcie dystrybucyjnym łączone
  w stos. Szacunkowa liczba przełączników 48-portowych: IDF2 -- 5 (240 portów
  na wymagane 203), IDF1 -- 3 (144 na 108), MDF -- 2 (96 na 77).
  Łącza w górę prowadzone są przez porty SFP, a łącza zagregowane
  rozłożone na różne przełączniki stosu.
\end{enumerate}

\paragraph{Łącza szkieletowe.}
\begin{enumerate}[resume]
  \item Łącze IDF2--MDF (światłowód wielomodowy, 204~m między budynkami):
  wymagane 194~Mb/s, po wzroście o 20\,\% 232,8~Mb/s. Przyjęto dwa łącza
  1000BASE-SX zagregowane (LACP), co daje 2~Gb/s i obciążenie poniżej
  12\,\%. Wymaga to dwóch par włókien w kablu międzybudynkowym (do
  potwierdzenia). Standard 1000BASE-SX działa na każdym typie włókna
  wielomodowego (OM1 do 275~m, OM2 do 550~m). Rzeczywista długość toru
  z IDF2 do MDF, z pionami w budynkach, będzie większa niż 204~m, więc przy
  włóknie OM1 zapas jest niewielki i trzeba zmierzyć tor. Jeśli inwentaryzacja
  potwierdzi włókno OM3/OM4, można zastosować 10GBASE-SR (do 300/400~m), który
  na OM1/OM2 nie osiąga tej odległości.
  \item Łącze IDF1--MDF (skrętka kat.~6, do 90~m): wymagane 99,5~Mb/s.
  Przyjęto dwa łącza 1000BASE-T zagregowane (LACP).
  \item Agregacja łączy zapewnia jednocześnie nadmiarowość: awaria jednego
  portu, wkładki lub pary włókien nie przerywa pracy budynku. Nie chroni
  natomiast przed przecięciem całego kabla, bo obie pary biegną tą samą trasą.
\end{enumerate}

\paragraph{Serwery.}
\begin{enumerate}[resume]
  \item Serwery lokalne w MDF, podłączone do przełącznika rdzeniowego
  łączami 2$\times$1~Gb/s (LACP). Najbardziej obciążony Serwer1 wymaga
  ok.\ 200~Mb/s.
  \item Serwery WWW i FTP w wydzielonej strefie DMZ, za zaporą sieciową.
\end{enumerate}

\paragraph{Podział logiczny (VLAN).}
\begin{enumerate}[resume]
  \item Osobne sieci VLAN dla: Drukarni, Korektorów, Konsultantów, sieci
  bezprzewodowej, drukarek, serwerów lokalnych, DMZ oraz zarządzania
  urządzeniami.
  \item Routing między VLAN-ami na przełączniku rdzeniowym w MDF, z listami
  kontroli dostępu (np.\ Serwer2 dostępny tylko dla Korektorów).
  \item Priorytetyzacja ruchu (QoS, znakowanie DSCP) dla VoIP
  i wideokonferencji na przełącznikach i zaporze.
  \item Dynamiczny przydział adresów (DHCP) dla stanowisk i urządzeń
  bezprzewodowych, adresy statyczne dla serwerów, drukarek i urządzeń
  sieciowych.
\end{enumerate}

\paragraph{Łącza do Internetu.}
\begin{enumerate}[resume]
  \item Łącze podstawowe: symetryczne, światłowodowe, o przepływności
  nominalnej co najmniej 100/100~Mb/s (wymagane 67,2/43,9~Mb/s, po wzroście
  o 20\,\% 80,4/51,7~Mb/s) i przepływności gwarantowanej (SLA) co najmniej
  40~Mb/s w obu kierunkach (wymagane 26,9/17,5~Mb/s, po wzroście
  32,2/20,7~Mb/s). Łącze musi być symetryczne, bo ruch w górę (43,9~Mb/s)
  przekracza wysyłanie typowych łączy asymetrycznych.
  \item Łącze zapasowe: od innego operatora, inną drogą fizyczną,
  o przepływności nominalnej co najmniej 100/100~Mb/s, bez gwarancji. Po
  awarii łącza podstawowego przejmuje cały ruch.
  \item Pula stałych publicznych adresów IPv4 (co najmniej /29) od operatora
  łącza podstawowego, dla serwerów WWW i FTP oraz translacji NAT. Od operatora
  łącza zapasowego co najmniej jeden stały adres publiczny. Po przełączeniu
  na łącze zapasowe serwery WWW i FTP pozostają dostępne dzięki zmianie
  rekordów DNS (krótki czas TTL).
\end{enumerate}

\paragraph{Bezpieczeństwo.}
\begin{enumerate}[resume]
  \item Zapora sieciowa nowej generacji (NGFW) na styku z Internetem: NAT,
  strefa DMZ, filtrowanie ruchu, system IPS, ochrona antywirusowa ruchu.
  \item Sieć bezprzewodowa w osobnym VLAN-ie, szyfrowanie WPA2/WPA3-Enterprise
  z uwierzytelnianiem 802.1X przez serwer RADIUS.
  \item Na przełącznikach: wyłączenie nieużywanych portów, port security,
  DHCP snooping, ochrona przed pętlami (RSTP, BPDU Guard).
  \item Oprogramowanie antywirusowe na stanowiskach i serwerach, regularne
  kopie zapasowe danych Serwera1 i Serwera2.
\end{enumerate}

\paragraph{Niezawodność i zasilanie.}
\begin{enumerate}[resume]
  \item Zasilacze awaryjne UPS w każdej szafie (MDF, IDF1, IDF2) oraz dla
  serwerów.
  \item Automatyczne przełączanie na łącze zapasowe przy awarii łącza
  podstawowego.
  \item Rezerwa portów i przepływności pozwalająca na wzrost liczby
  użytkowników o 20\,\% (tabela~\ref{tab:wzrost}).
  \item Przełącznik rdzeniowy i zapora z redundantnymi zasilaczami.
  Pojedynczy przełącznik rdzeniowy i pojedyncza zapora pozostają punktami
  awarii, ocena tego ryzyka i ewentualne zdublowanie tych urządzeń zostaną
  przedstawione w analizie bezpieczeństwa i niezawodności (punkt~5.6).
\end{enumerate}

\section{Projekt sieci}

\subsection{Projekt logiczny sieci wraz z opisem koncepcji rozwiązania i uzasadnieniem}
% Etap 3

\subsection{Wybór urządzeń sieciowych}
% Etap 3

\subsection{Projekt adresacji IP}
% Etap 3

\subsection{Projekt konfiguracji urządzeń}
% Etap 4

\subsection{Projekt podłączenia do Internetu}
% Etap 4

\subsection{Analiza bezpieczeństwa i niezawodności sieci}
% Etap 4

\subsection{Kosztorys}
% Etap 4

\section{Karty katalogowe proponowanych urządzeń}
% Etap 4

\end{document}
